% sec-08-limitations-future-work.tex - Section 8.  FINAL stage.
% Figure 25 drawn; Table 25 populated.  Senior-author pass: the two smaller
% limits are moved up into the Limitations subsection where they belong rather
% than trailing the Future Work subsection, and the register of seven absent
% instruments is expanded so \S8 and \S6 agree.
\section{Limitations and Future Work}
\label{sec:limitations}

\subsection{Limitations}

Two limitations are structural rather than incidental. The first is that a
single model family, Claude Code Opus 5, wrote every line of source in this
paper and in the eight-stage build that produced it. A defect in that model's
idiom appears in every artifact at once, and no amount of internal consistency
checking inside one family detects a shared blind spot. The source study for \S7
records the same class of limit from the other direction: its first working
notebook used a model whose parameters were ``nearly 6 years'' old and a newer
alternative ``may have been a better replacement'', which the author attributes
to the generating model relying ``on older information from pre-trained
parameters as opposed to a more exhaustive online model search process''
\cite{aiprgliomatc}. It also records a scale ceiling: ``there was an upper limit
when working with over 1500 lines of code'', and truncation appeared ``when
including additional documents beyond the notebook for consideration'', so that
``the most ideal responses were typically received when the author defined the
problem the most directly with the least number of file attachments''
\cite{aiprgliomatc}. The mitigation is the two independent reviewer
manufacturers of \S7, and the residual is the case where all three agree and all
three are wrong.

The second limitation is text watermarking. Generated text may carry a
provenance watermark, detection policy is not settled, and journals differ on
what disclosure they require. This paper's answer is not to argue about
detection. It is to disclose the model on the cover page of every deposited
work, to make the commit history rather than a claim of authorship the primary
record, and to deposit to a repository and a preprint rather than through a
journal gate. Figure 25 pairs each exposure with the mitigation that reduces it,
and Table 25 states the same register in citable text with the residual that
survives each.

Two further limits deserve statement even though they are smaller. This paper
reports on documents rather than on patients: no participant has been enrolled,
no operation has been performed under these protocols, and every clinical number
in \S\ref{sec:ind} and \S\ref{sec:protocol} is either a published external
result or a planned design parameter. And the review regime of
\S\ref{sec:peerreview} was measured on a glioblastoma machine learning workflow
rather than on a regulatory dossier; the method transfers, and the measurement
does not. Both limits are removable only by doing the trial, which is what the
funding is for.

% ---------------------------------------------------------------------------
% FIGURE 25.  diagrams (python)-type, clustered infrastructure.  Two exposure
% clusters paired to two mitigation clusters across six shared y values, so
% every pairing edge is horizontal, converging on one residual node.
% ---------------------------------------------------------------------------
\begin{appfloat}[!tb]
\begin{appfig}[diagrams (python)-type / clustered infrastructure]
\dgnodeg{dgtilek}{1.05}{1.55}{\glyphai}{One model family writes all source}{a1}
\dgnodeg{dgtilek}{1.05}{-0.35}{\glyphdoc}{One context window bounds each stage}{a2}
\dgnodeg{dgtilek}{1.05}{-2.25}{\glyphcloud}{Vendor availability is a project risk}{a3}
\begin{scope}[on background layer]
\node[dgcluster2,fit=(a1)(a1l)(a2)(a2l)(a3)(a3l)] (ce1) {};\end{scope}
\node[dgctitle2,anchor=south west] at ([yshift=1.5mm]ce1.north west) {Exposure 1, single-vendor production};
\dgnode{dgtilem}{8.35}{1.55}{\glyphteam}{Two independent reviewer makers}{m1}
\dgnode{dgtilem}{8.35}{-0.35}{\glyphlink}{Machine-readable, portable specifications}{m2}
\dgnode{dgtilem}{8.35}{-2.25}{\glyphlock}{Every artifact deposited under a DOI}{m3}
\begin{scope}[on background layer]
\node[dgcluster,fit=(m1)(m1l)(m2)(m2l)(m3)(m3l)] (cm1) {};\end{scope}
\node[dgctitle,anchor=south west] at ([yshift=1.5mm]cm1.north west) {Mitigation 1};
\draw[trialslate,dashed,line width=0.4pt] (-0.35,-3.40) -- (14.80,-3.40);
\dgnodeg{dgtilek}{1.05}{-4.55}{\glyphsignal}{Generated text may carry a watermark}{b1}
\dgnodeg{dgtilek}{1.05}{-6.45}{\glyphmon}{Detection policy is not yet settled}{b2}
\dgnodeg{dgtilek}{1.05}{-8.35}{\glyphbank}{Journals differ on disclosure}{b3}
\begin{scope}[on background layer]
\node[dgcluster2,fit=(b1)(b1l)(b2)(b2l)(b3)(b3l)] (ce2) {};\end{scope}
\node[dgctitle2,anchor=south west] at ([yshift=1.5mm]ce2.north west) {Exposure 2, text watermarking};
\dgnode{dgtilem}{8.35}{-4.55}{\glyphdoc}{Provenance disclosed on the cover page}{n1}
\dgnode{dgtilem}{8.35}{-6.45}{\glyphdb}{Commit history is the primary record}{n2}
\dgnode{dgtilem}{8.35}{-8.35}{\glyphserver}{Preprint and repository, not a journal gate}{n3}
\begin{scope}[on background layer]
\node[dgcluster,fit=(n1)(n1l)(n2)(n2l)(n3)(n3l)] (cm2) {};\end{scope}
\node[dgctitle,anchor=south west] at ([yshift=1.5mm]cm2.north west) {Mitigation 2};
\dgnodew{dgtiled}{12.95}{-4.45}{\glyphchart}{Residual: funding decision risk}{res}
\draw[trialburg,line width=0.6pt] (res.center) circle (7.5mm);
\foreach \a/\b in {a1/m1,a2/m2,a3/m3,b1/n1,b2/n2,b3/n3}{
  \draw[dgedge] (\a.east) -- node[above,font=\tiny\sffamily,fill=trialwhite,inner sep=1pt] {reduced by} (\b.west);}
\draw[dgedged] (m3.east) -- (10.05,-2.25) -- (res.north west);
\draw[dgedged] (n3.east) -- (10.05,-8.35) -- (res.south west);
\pnote{0.35}{-9.55}{Exposures and mitigations share six y values exactly, so every pairing edge is
horizontal and none is oblique.\\After six mitigations exactly one consequence
remains, and it is a funding decision risk rather than a safety one.}
\end{appfig}
\vspace{-0.6cm}
\figcaption{Figure 25. Two structural exposures, single-vendor production and
text\\watermarking, each drawn with the mitigation that reduces it and its
residual.}
\end{appfloat}

\begin{apptable}
\begin{tabularx}{\textwidth}{>{\raggedright\arraybackslash}p{4.45cm} >{\raggedright\arraybackslash}p{4.55cm} >{\raggedright\arraybackslash}X}
\headrow \thead{Exposure} & \thead{Mitigation} & \thead{Residual after mitigation}\\
\midrule
One model family writes all source & Two independent reviewer manufacturers at defined milestones & A defect all three share would survive \\
One context window bounds each stage & Machine-readable, portable specifications & Re-execution costs a stage, not a project \\
Vendor availability is a project risk & Every artifact deposited under a DOI as produced & Deposited work is unaffected \\
Generated text may carry a watermark & Provenance disclosed on every cover page & Disclosure does not remove the watermark \\
Detection policy is not settled & Commit history is the primary record & Policy may change after deposit \\
Journals differ on disclosure & Deposit to preprint and repository & A funder requiring journal review is not reached \\
\end{tabularx}
\end{apptable}
\vspace{-0.6cm}
\tabcap{Table 25. Six exposures, the mitigation applied to each, and the residual\\that
survives it, ending in one funding decision risk.}

\subsection{Future work}

Future work is one thing: securing funding for the PDAC hybrid clinical trial
described in \S3 and \S4. The documents exist. The IND is written, both
protocols are written, the statutory framework is drafted, and the applications
are filed. What is absent is money and the institutional instruments money buys.
The capitalization plan states the shape of the ask precisely: five Phase I
milestones summing to \$306,000 over nine months, then seven Phase II milestones
summing to \$1,300,000 over twenty-four, each milestone closed by three
concurrent evidence branches that must join before an audit comparison
\cite{capplan2026}. Seven instruments are absent, and not one can be obtained by
the company alone: a qualified site, an operating room, a pathology service, an
institutional review board of record, a drug supply agreement, a device
acceptance certificate, and a clinical trial agreement.

No further document production changes that. This paper is therefore the last
document in the sequence rather than the next one in it, and the sequence ends
here on purpose.

The seven instruments named above are the register the whole paper converges on.
Four of them, the site, the operating room, the pathology service and the
institutional review board of record, are the ones \S\ref{sec:ind} and
\S\ref{sec:protocol} identified independently; three more, the drug supply
agreement, the device acceptance certificate and the clinical trial agreement,
are contractual rather than clinical. None can be produced by writing, which is
the sense in which this paper is the end of a documentary sequence rather than a
step in one. The method described in \S\ref{sec:methods} solved the problem it
was built for, and the problem that remains is of a different kind.
